\chapter{La transcription automatique, une solution privilégiée pour implémenter la base de données ?}

Dans notre utilisation de l'IA pour traiter les factures Rondel, la principale technologie expérimentée est la reconnaissance d'écritures manuscrites, ce qui semple plutôt logique au vu de la nature du corpus. S'il est essentiel de revenir sur l'histoire de cette technologie, ce chapitre vise surtout à présenter les différents tests réalisés sur les factures Rondel et ce qui a été mis en place par le département pour poursuivre le projet. 

\section{La \emph{Handwritten Text Recognition}, ou la reconnaissance automatique d'écritures manuscrites}

La grande majorité des factures traitées jusqu'à présent est de nature manuscrite. Pour les transcrire, il convient donc de s'intéresser à une technologie nommée \emph{Handwritten Text Recognition} (HTR), utilisant l'intelligence artificiele et spécialisée dans la reconnaissance et la transcription d'écritures manuscrites. Cette technologie est à distinguer de l'\emph{optical character recognition} (OCR), qui est plus ancienne\footcite[Voir][La naissance de l'OCR est souvent attribuée au physicien Emmanuel Goldberg qui, durant la Première Guerre mondiale, invente une machine permettant de lire des caractères et de les convertir en code télégraphique]{lefevre_locr_2023} et vise à transcrire des caractères imprimés\footcite[p. 1]{gutteridge_judge_2026}. La HTR est initialement fondée sur une collaboration entre l'IA et l'être humain, qui a pour rôle d'améliorer les performances de la machine par l'utilisation de corpus d'entraînement\footcite{pinche_2022}. Tout comme l'OCR, les outils traditionnels de la HTR fonctionnent en trois étapes. Cela commence par un prétraitement de documents numérisés, suivi d'une analyse et d'une segmentation de leur mise en page pour enfin en transcrire le texte\footcite[p. 1]{kim_early_2025}. 
Cependant, les évolutions actuelles dans ce domaine montrent que tous les outils de HTR connus jusqu'à présent sont aujourd'hui largement surpassés par les modèles de langage\footnote{Deux principaux types de modèle de langage sont à distinguer : les LLM (Large Language Modele), qui ne traitent que des données textuelles et les VLM (Vision Language Modele) qui, associant les technologies des LLM à des outils de vision, sont aussi capables de traiter des images.}, modèles d'IA générative permettant notamment de transcrire, de classer et de comprendre des documents\footcite[p. 1]{crosilla_benchmarking_2025}. Le type de modèle de langage ayant particulièrement fait évoluer la transcription automatique est le VLM, par sa capacité à traiter des images, ce que ne peut pas réaliser le LLM. Ces modèles n'ont pas besoin d'être entraînés et ne se contentent pas de transcrire l'écriture manuscrite. Ils sont également capables d'identifier la mise en page des documents, d'extraire, de structurer et d'analyser le sens des informations et de les exporter\footcite{sebastien_cretin_2025}. De plus, tout le travail n'est réalisé qu'en une seule étape du point de vue utilisateur\footcite[p. 1]{kim_early_2025}. Il faut surtout retenir que les capacités de ces modèles de langage s'améliorent très vite. En effet, la littérature scientifique sur le sujet affirmait encore en octobre 2024 que les performances des modèles de langage, \emph{Gemini}\index{Gemini} en particulier, étaient nettement inférieures à celles des modèles entraînés pré-existants, mais qu'elles étaient prometteuses\footcite[p. 4-6]{li_handwriting_2024}. L'outil semble s'améliorer constamment depuis, au point d'être actuellement comparable, selon certaines sources, à des performances humaines\footcite{humphries_gemini_2025}. 


Ainsi, revenir sur l'histoire et les évolutions de la HTR met en évidence le fait que les grands modèles de langage, en particulier les VLM, sont aujourd'hui une solution simple et rapide à utiliser, pouvant réaliser de nombreuses tâches et de plus en plus performante. 


\section{Les différents outils de reconnaissance automatique expérimentés pour traiter les factures Rondel} 

Quelque soit les solutions expérimentées, envisager un traitement informatique de notre corpus suppose d'identifier précisément nos besoins en matière d'automatisation. L'exploitation informatique des factures est surtout utile pour alimenter la table\enquote{ligne de facture} de la base, qui comporte le plus d'informations issues des factures et qui vise à identifier les oeuvres. Ainsi, nos attendus vis-à-vis de l'IA sont l'analyse de la mise en page des factures, afin de comprendre et de transcrire un maximum d'informations susceptibles de remplir cette table. À partir de ce travail, nous souhaitons que l'IA structure les informations dans un format permettant de les importer directement dans la base. Comme évoqué précédemment, il est préférable que les informations issues de cette transcription soient structurées au format CSV, manière la plus simple d'importer des données en masse dans \emph{Heurist}\index{Heurist}. Pour que les imports soient possibles, nous avions d'autres besoins plus précis, à savoir de corriger et d'adapter certaines données pour qu'elles puissent être insérées dans la base. L'attendu majeur est que l'IA lie chaque ligne de facture à l'identifiant \emph{Heurist}
(\emph{Heurist-ID})\footnote{Voir chapitre I, partie II, \enquote{\nameref{types donnees}}.} de l'enregistrement \enquote{facture} créé au préalable, à laquelle la ligne doit être liée. Cela permet de lier directement les lignes à leur facture lors de l'import du fichier CSV\footnote{Voir page 28 à 31 de \enquote{\nameref{documentation}}.}.


La nécessité d'un outil utilisable rapidement et en mesure de répondre simultanément à plusieurs besoins a fait des modèles de langage la technologie priorisée dans l'exploitation des factures Rondel. En effet, entraîner un modèle de HTR n'est pas le but de notre projet et était trop long pour être réalisé sur le temps de mon stage. Le principal modèle de langage utilisé pour exploiter notre corpus est Gemini\index{Gemini}. Sur la durée de mon stage, plusieurs de ses modèles ont pu être expérimentés, tels que \emph{Gemini}\index{Gemini} 3.5 Flash, sorti en mai 2026, moment où nous avons commencé à utiliser l'IA. Les derniers modèles de \emph{Gemini}\index{Gemini} sont capables de réaliser des tâches parfaitement adaptées à l'exploitation des factures Rondel, à savoir la compréhension d'images, l'extraction et la classification d'informations\footcite{noauthor_liste_nodate}.
Les possibilités de traiter les factures Rondel grâce à \emph{Gemini}\index{Gemini} ont été explorées par différents prompts, incluant progressivement des consignes de plus en plus précises répondant aux exigences de notre base. Le prompt le plus complet vise à remplir le plus de champs possibles de la table \enquote{ligne de facture}. 
Ce prompt est visible en annexe, accompagné d'un exemple de facture transcrite, et du fichier CSV extrait de cette transcription\footnote{Voir \enquote{\nameref{prompt gemini}}.}. La facture, datant du 22 juin 1920, est produite par la librairie Hermal, qui a vendu beaucoup d'oeuvres à Auguste Rondel. La plupart des tests d'IA ont été réalisés sur les factures produites par Hermal, car les descriptions des oeuvres sont très complètes, ce qui donne un bon aperçu des capacités de \emph{Gemini}\index{Gemini}. L'exemple montré en annexe illustre la finesse de l'analyse textuelle faite par \emph{Gemini}\index{Gemini}. L'outil a parfaitement transcrit toutes les informations sur chaque oeuvre et les a systématiquemet rangées dans les bonnes colonnes. \emph{Gemini}\index{Gemini} parvient même à détecter certaines subtilités. C'est le cas de la distinction entre auteur et éditeur, ce qui transparaît en particulier dans la première ligne, où seul un éditeur est mentionné. On voit aussi dans les cinquième et sixième ligne que Gemini sait faire la distinction entre la date d'édition de l'oeuvre et une date faisant partie du titre de l'oeuvre. Comme Gemini ne parvient pas à interroger le catalogue général de la BnF\footnote{Ce point sera détaillé dans le chapitre suivant.}, le prompt tente malgré tout de relier chaque oeuvre à sa sous-série d'appartenance dans la collection par un moyen détourné. Il demande à \emph{Gemini}\index{Gemini} de relier chaque oeuvre à la sous-série de la collection Rondel à laquelle elle est la plus susceptible d'appartenir, ce qu'on peut voir dans la colonne \enquote{Série/sous-série} du fichier CSV. Même s'il ne s'agit que d'hypothèses, cette méthode fonctionne plutôt bien. Après vérification, quatre des six oeuvres apparaissant dans la facture du 22 juin 1920 ont été reliées à la bonne sous-série. Cependant, comme pour tous les autres traitements demandés à \emph{Gemini}\index{Gemini}, une vérification humaine est indispensable avant le moindre import de ces données dans la base, même si l'outil est très performant. Ce prompt, qui a prouvé son efficacité dans la transcription d'une facture afin d'alimenter la table \enquote{ligne de facture}, a donc été retenu et expérimenté plusieurs fois. Des consignes supplémentaires ont également été testées et peuvent être ajoutées pour améliorer le prompt. On peut par exemple demander à \emph{Gemini}\index{Gemini} de ne pas transcrire certaines lignes, si elles ne font pas référence à des oeuvres. \emph{Gemini}\index{Gemini} est également capable de transcrire plusieurs factures à la fois, si on lui indique correctement quels identifiants \emph{Heurist} écrire sur quelles lignes, pour que les lignes soient associées aux bonnes factures une fois dans la base. Il faut donc remplacer l'instruction liée à la colonne \enquote{H-ID facture} comme dans l'exemple suivant : 

\begin{quote}
    Dans la colonne "H-ID facture", ajoute 1718  sur toutes les lignes pour les factures datant du 12 novembre 1919, 1719 sur toutes les lignes pour celles datant du 4 novembre et 1720  sur celles datant du 13 octobre. 
\end{quote}

Ainsi, \emph{Gemini}\index{Gemini} est de loin l'outil le plus performant, rapide et simple d'utilisation essayé pour transcrire les factures Rondel. Malgré cela, nous nous sommes aussi intéressés aux alternatives \emph{Open Source} existantes pour réaliser ces mêmes tâches. Le principal outil expérimenté est \emph{NuExtract3}\index{NuExtract3}, publié le 19 mai 2026 par \emph{NuMind}. Il s'agit d'un modèle VLM spécialisé dans l'extraction de documents. Il réunit l'extraction structurée et la reconnaissance de caractères\footcite{numind_nuextract_nodate}. Cet outil a été testé via \emph{Hugging Face}\index{Hugging Face}, plateforme dédiée au partage de jeux de données, d'applications et modèles d'apprentissage par intelligence artificielle\footcite{noauthor_nuextract_nodate}. Cette plateforme permet une prise en main relativement simple : il suffit d'entrer le document que l'on souhaite transcrire dans la catégorie \emph{input}, de saisir au format JSON la structure des informations que l'on souhaite extraire dans \emph{Schema \& Instructions} et de cliquer sur \emph{Extract JSON}. Nous avons testé cet outil sur des factures produites par différents libraires, mais les tests les plus convaincants sont ceux réalisés sur celles de la librairie florentine \emph{de Marinis}, ou sur celles de la librairie de l'Ancien temps, située à Paris. Cela s'explique par le peu d'informations contenues dans ces factures et le fait qu'elles se présentent souvent sous forme tabulaire, avec une délimitation claire entre les différents types d'informations. Même s'il reste quelques fautes à corriger avant d'importer dans la base, l'outil est plutôt efficace. Un exemple de test réalisé sur une facture, contenant la facture elle-même, le schéma JSON utilisé pour extraire et le résultat est disponible en annexe\footnote{Voir dans l'annexe E, I.2 : \enquote{\nameref{nuextract3}}.}. Le résultat de l'extraction est structuré sous la forme d'un schéma JSON, constitué de clés (les noms des champs de la table) et de valeurs (les valeurs de ces champs pour chaque enregistrement). L'import des données JSON directement dans la base \emph{Heurist} est possible, mais cela était difficile à réaliser correctement. L'alternative choisie pour avoir des données prêtes à être importées dans la base a donc consisté en un script Python permettant de convertir le schéma JSON en un fichier CSV, lui aussi détaillé en annexe. Malgré tout, \emph{NuExtract3}\index{NuExtract3} n'a été efficace que sur certaines factures parfaitement structurées et pas trop denses. Sur d'autres factures, les transcriptions faites par cet outil étaient beaucoup moins précises, contenaient des erreurs et étaient mal structurées. Cet outil n'est mentionné ici que parce qu'il est plus performant que les autres outils \emph{Open Source} testés sur les factures Rondel. Malgré tout, sa pertinence reste nettement inférieure à celle de \emph{Gemini}\index{Gemini}, car il est plus lent, transcrit moins bien et n'est pas conçu pour réaliser automatiquement tous les traitements nécessaires.  

Toute transcription automatique des factures est suivie d'un import, de préférence au format CSV, dans la base. Pour que l'import soit réussi, il est très important de sélectionner le bon séparateur, caractère permettant de diviser le texte entre les différentes colonnes, ainsi que le bon encodage, qui assure que tous les caractères soient correctement lus par le logiciel. Une fonction de prévisualisation du futur import, appelée \emph{Analyse data}, permet de s'assurer que le fichier s'importera correctement dans le logiciel. Comme évoqué au début de cette partie, les étapes de cet import sont détaillées dans la documentation sur les usages ultérieurs de la base\footnote{Voir page 28 à 31 de \enquote{\nameref{documentation}}}. 


Tous ces essais d'utilisation de l'IA dans le traitement des factures Rondel ont permis de dresser un constat principal. Celui-ci est, nous l'avons vu, que \emph{Gemini}\index{Gemini} surpasse de loin tous les autres outils essayés dans le cadre de ce projet, dans toutes les phases du traitement des données. Or, il convient de rester critique envers cet outil, car celui-ci présente certaines limites qui n'ont rien à voir avec ses performances. Premièrement, \emph{Gemini}\index{Gemini} est très consommateur en énergie. Même s'il n'existe pas de données chiffrées publiques concernant cette consommation, les estimations placent \emph{Gemini}\index{Gemini} comme l'un des outils les plus consommateurs au monde. Le site ComparIA, géré par le gouvernement français et permettant de comparer les IA conversationnelles, a classé la taille du modèle en \emph{XL}\footcite{noauthor_classement_nodate}. La taille d'un modèle d'IA influe nécessairement sur sa consommation énergétique : plus un modèle est lourd, plus il consomme. De plus, \emph{Gemini}\index{Gemini} n'est pas une solution \emph{Open Source} et Google en est complètement propriétaire, ce qui explique notamment le manque de transparence au sujet de sa consommation. Il faut donc avoir conscience, lorsqu'on utilise \emph{Gemini}\index{Gemini}, que toutes les données traitées sont transmises à Google. Les modèles \emph{Open Source}, quant à eux, n'engendrent pas de fuites de données et consomment beaucoup moins que \emph{Gemini}\index{Gemini}. Par exemple, la taille de \emph{NuExtract3}\index{NuExtract3} est estimée à quatre milliards de paramètres\footcite{numind_nuextract_nodate}.
En revanche, ils sont nettement moins performants sur le plan technique, comme nous l'avons vu. 

Nos expérimentations elles-mêmes ont aussi des limites à prendre en compte. Même s'il est très important, le traitement par IA de notre corpus n'est pas l'objet principal de ce projet, ce qui explique l'absence de statistiques sur l'efficacité des différents outils essayés. Cela aurait nécessité un travail dit de \enquote{vérité terrain}, passant par la sélection d'un échantillon représentatif du corpus, et la comparaison de sa transcription manuelle et de sa transcription automatique. Nous avons simplement réalisé une preuve de concept, c'est-à-dire une démonstration visant à prouver la faisabilité d'une idée\footcite{definition_PoC}. En somme, nos analyse sont qualitatives, et non quantitatives. 



\section{La numérisation : une étape nécessaire dans le traitement des factures par intelligence artificielle} 

Un potentiel traitement massif de ces factures de libraires par IA nécessite de mettre en place leur numérisation. Au début de ce projet, aucune facture du corpus n'était numérisée. La plupart des tests d'IA ont donc été réalisés sur des photographies de factures, ce qui ne peut pas être une solution à long terme. Nous avons demandé et obtenu la possibilité de numériser les 2000 premières factures du corpus, c'est-à-dire les cinq premières boîtes. Avant d'être numérisées, les factures de chaque dossier ont été classées par ordre chronologique  lorsqu'elles ne l'étaient pas déjà, pour qu'elles apparaissent dans cet ordre sur Gallica\index{Gallica}. 
Quelques consignes supplémentaires ont été données. Par exemple, nous avons choisi de ne pas numériser les mandats de paiement, ni les documents autres que des factures contenus dans les dossiers (accusés de réception, avis de débit, ...). Ont évidemment été exclus de cette numérisation les dossiers consacrés à d'autres commerces que des librairies, qui ne seraient d'aucune utilité au remplissage de la base (ateliers de relieurs, imprimeries, ...). Ces 2000 factures de libraires sont dans Gallica\index{Gallica} et sont facilement accessibles depuis le catalogue BnF Archives et manuscrits\index{Bnf Archives et manuscrits}. Un exemple de dossier numérisé est disponible en annexe\footnote{Voir annexe E, II. \enquote{\nameref{numerisation}}.}. Ce travail de numérisation a pour objectif d'être poursuivi sur le reste du corpus dans le cadre d'une continuation de ce projet. Malgré le caractère expérimental de l'usage de l'IA dans l'exploitation des factures Rondel, cette première numérisation ouvre la voie, par les solutions actuelles ou futures, à la pérennisation de ces procédés. 



\vspace{1,5cm}
En somme, les tentatives d'utilisation de l'IA pour exploiter les factures Rondel se sont d'abord appuyées sur un état de l'art, afin d'être en phase avec les pratiques dans ce domaine. Celles-ci ont permis d'évaluer quels sont les outils adaptés aux traitements dont on a besoin, et les contraintes liées aux différents modèles. Cependant, nos essais sont limités pour l'instant et demandent à être déployés à une échelle plus large. 